%% kapa: exam=1 round=21 year=2010 period=2 form=A subject=appraisal_law subject_ko=부동산관계법규 seq=21 number=61
%% source: Q-Net 감정평가사 기출문제 (cst00309) artlSeq=5024181 file=extracted/1008/5024181/2010년 제21회 감정평가사 1차시험 기출문제/2교시.hwp
%% conversion: kordoc 4.24.51 -> tools/gen1.py (auto); review=proofread-gpt-sol; flags=none
\begin{question}{61}
부동산 가격공시 및 감정평가에 관한 법령상 개별공시지가를 결정ㆍ공시함에 있어서 감정평가업자의 검증에 관한 설명으로 옳지 않은 것은?
\begin{choices}
  \item 시장ㆍ군수 또는 구청장은 개별공시지가를 결정․공시하기 위하여 검증을 의뢰하고자 할 때에는 전체 개별토지의 가격에 대한 지가현황도면 및 지가조사자료를 제공하여야 한다.
  \item 검증의뢰를 받은 감정평가업자는 비교표준지의 선정에 관한 사항 등을 검토ㆍ확인하고 의견을 제시하여야 한다.
  \item 시장ㆍ군수 또는 구청장은 감정평가업자의 검증이 필요없다고 인정되는 때에는 지가의 변동상황 등 대통령령이 정하는 바에 따라 감정평가업자의 검증을 생략할 수 있다.
  \item 개발사업이 시행되거나 용도지역ㆍ지구가 변경되는 등의 사유가 발생한 토지에 대하여는 감정평가업자의 검증을 생략할 수 있다.
  \item 시장ㆍ군수 또는 구청장은 당해 지역의 표준지의 공시지가를 조사․평가한 감정평가업자 또는 대통령령이 정하는 감정평가실적 등이 우수한 감정평가업자에게 의뢰하여야 한다.
\end{choices}
\end{question}
\answer{4}
